%% 自動生成（gen.py）
\bktocchap{7-0}{表のままで計算できたら？}{九九の表は、行列の入口}{5}{3}{3}{8}{J}
\bktocpart{第1部}{行列の計算}{22}
\bktocchap{7-1}{数の表で計算する：行列}{買い物の合計を、表のかけ算で一度に求める}{5}{3}{3}{24}{C}
\bktocchap{7-2}{行列×ベクトル：形を変える機械}{4つの数で、絵をのばしたり回したりできる}{5}{4}{4}{64}{C}
\bktocchap{7-3}{行列×行列：変換をつなぐ}{「回してのばす」と「のばして回す」はちがう}{5}{4}{4}{102}{C}
\bktocchap{7-4}{行列式：面積が何倍になるか}{0になると、平らにつぶれて元にもどせない}{5}{4}{4}{138}{C}
\bktocchap{7-5}{逆行列で連立方程式を解く}{式が100個あっても、コンピュータの手順は同じ}{5}{4}{4}{186}{C}
\bktocpart{第2部}{回転と座標の変換}{236}
\bktocchap{7-6}{回転行列：くるっと回す}{30°回して45°回すと、75°回したのと同じ}{4}{4}{4}{238}{C}
\bktocchap{7-7}{3次元の回転と座標の変換}{ドローンの向きを、3つの回転で表す}{6}{5}{5}{274}{C}
\bktocpart{第3部}{固有値と行列の使いみち}{310}
\bktocchap{7-8}{固有値と固有ベクトル：向きの変わらない方向}{変換しても、向きが変わらない矢印がある}{5}{5}{5}{312}{C}
\bktocchap{7-9}{対角化：まっすぐな方向に分けて考える}{こみいった変換も、のびるだけの変換に分けられる}{5}{5}{5}{354}{C}
\bktocchap{7-10}{行列のくり返し：漸化式とオイラー法を行列で}{何回もかけるとどうなるかは、固有値が教えてくれる}{5}{5}{5}{386}{C}
\bktocchap{7-11}{行列の使いみち}{ロボット、電気回路、画像…行列はどこで働いている？}{5}{3}{5}{430}{R}
\par\vspace{2mm}
\bktocbacklist{\bkbi{こたえ}{450}\bkbi{\R{記号}{き|ごう}カード}{464}\bkbi{よく\R{使}{つか}う\R{定数}{てい|すう}・\R{単位}{たん|い}・\R{公式}{こう|しき}}{469}\bkbi{できるようになったことチェック\R{表}{ひょう}}{484}\bkbi{つぎの\R{巻}{かん}へ}{486}\bkbi{さくいん}{488}}
